% Run pdfLaTeX from this directory to resolve the shared inputs.
\makeatletter
\def\input@path{{../../../preamble/}}
\makeatother
\input{preamble.tex}
\input{graphs/highlights.tex}
% \GTenHighlightsfalse % Uncomment to disable instructional highlights.

\geometry{letterpaper,margin=0.25in}
\pagestyle{plain}
\linespread{1.04}
\setlength{\parskip}{4pt}
\hypersetup{pdftitle={Grade 10 Term 2 Lesson 1 Classroom Activity: Outcomes, Support and PMFs}}
\newif\ifprintsolutions
\printsolutionstrue % Change to \printsolutionsfalse for the student handout.

\setlength{\tabcolsep}{6pt}
\setlength{\extrarowheight}{2pt}
\newcolumntype{Y}{>{\centering\arraybackslash}X}
\newcommand{\stageheading}[1]{%
    \par\addvspace{0.65\baselineskip}%
    {\centering\large\bfseries #1\par}\nobreak
}
\newcommand{\problemtitle}[3]{%
    \par\addvspace{0.4\baselineskip}%
    % Keep the original line height and depth while drawing the padding.
    {\setlength{\fboxsep}{2pt}%
     \raisebox{0pt}[\dimexpr\height-\fboxsep\relax][\dimexpr\depth-\fboxsep\relax]{%
         \colorbox{black!10}{\textbf{#1. #2}}}}%
    \hfill\textbf{Type #3}\par\nobreak
}
\newcommand{\taskheading}[1]{%
    \textbf{Task #1. \ifcase#1\or
    Identify the outcomes and their probabilities.\or
    Identify the random variable values and their probabilities.\or
    Find a specific probability.\fi}%
}
\newcommand{\task}[2]{\par\noindent\taskheading{#1} #2\par}
\newcommand{\responseline}{\noindent\rule{0pt}{1.55em}\hrulefill\par}
\newcommand{\responselines}[1]{%
    \par\begingroup\setlength{\parskip}{0pt}%
    \responseline\ifnum#1=2\responseline\fi
    \endgroup
}
\newcommand{\nextpracticeproblem}{%
    \par\vspace{1.6in}\vfill\hrule\medskip
}

% Keep the conditional assignment inside a macro so printsolutions can skip it.
\newcommand{\RestoreClassroomHighlights}{%
    \let\ifGTenHighlights\ClassroomHighlightPreference
}

\begin{document}
% Suppress student-work styling while retaining the master preference.
\let\ClassroomHighlightPreference\ifGTenHighlights
\GTenHighlightsfalse

{\centering
    {\Large\bfseries C1 -- Outcomes, Support and PMFs\par}
    Global Economics \textbar{} Grade 10 \textbar{} Term 2 \textbar{} Lesson 1\par
}

\textbf{C1:} Characterises discrete random variables in economic and financial contexts by identifying their support, probability mass function and possible outcomes.

A random variable assigns a number to each complete outcome. Its probability mass function (PMF) gives the probability of each value. The support contains the values with positive probability.

\stageheading{Worked examples}

	\problemtitle{1}{Changes to the till balance}{A}
	A shop sells a jacket for \(12\) USD. Exactly one of the events in the table happens by the end of the return period.

\noindent
\begin{minipage}[t]{0.48\linewidth}
	\task{1}{Read the four possible outcomes and their probabilities from the table.}
	
	\begin{tabularx}{\linewidth}{|l|Y|Y|}
		\hline
		Complete outcome & Probability & \(X\) (USD) \\ \hline
		Kept & \(0.37\) & \(0\) \\ \hline
		Even exchange & \(0.21\) & \(0\) \\ \hline
		Upgrade & \(0.16\) & \(3\) \\ \hline
		Full refund & \(0.26\) & \(-12\) \\ \hline
	\end{tabularx}
\end{minipage}
\hfill
\begin{minipage}[t]{0.48\linewidth}
	\task{2}{Map the outcomes and group the values. Let \(X\) be the change in the shop's till balance after the original payment, in USD. Keeping the jacket or exchanging it at the same price gives \(X=0\). The upgrade gives \(X=3\) and the full refund gives \(X=-12\). Add the probabilities of the two outcomes that give \(X=0\).}
	
	\begin{tabularx}{\linewidth}{|c|Y|Y|}
		\hline
		\(x\) (USD) & Outcomes giving \(x\) & Calculation and mass \(p_X(x)\) \\ \hline
		\(-12\) & Full refund & \(0.26\) \\ \hline
		\(0\) & Kept and Even exchange & \(0.37+0.21\)\newline \(=0.58\) \\ \hline
		\(3\) & Upgrade & \(0.16\) \\ \hline
	\end{tabularx}
	
	\medskip
	The support is \(\{-12,0,3\}\) in USD.
	All masses are nonnegative and
	\(0.26+0.58+0.16=1\).
\end{minipage}

\task{3}{Find the probability that the balance stays the same after the original payment. The PMF gives \(P(X=0)=p_X(0)=0.58\). The balance stays the same in \(58\%\) of these sales.}

\problemtitle{2}{Invoice payment adjustment}{B}
Two customers independently pay their invoices by the deadline (\(\mathrm{T}\)) with probabilities \(0.70\) and \(0.55\), respectively, or pay late (\(\mathrm{L}\)). The business earns an \(8\) USD bonus if both pay by the deadline, pays a \(6\) USD charge if both pay late, and has no adjustment otherwise.

\noindent
\begin{minipage}[t]{0.48\linewidth}
	\task{1}{List every pair in customer order and calculate its probability. Late-payment probabilities are \(1-0.70=0.30\) and \(1-0.55=0.45\). Independence lets us multiply within each pair.}
	
	\begin{tabularx}{\linewidth}{|l|Y|Y|}
		\hline
		Ordered outcome & Probability & \(Y\) (USD) \\ \hline
		\(\mathrm{TT}\) & \(0.70\cdot0.55\)\newline \(=0.385\) & \(8\) \\ \hline
		\(\mathrm{TL}\) & \(0.70\cdot0.45\)\newline \(=0.315\) & \(0\) \\ \hline
		\(\mathrm{LT}\) & \(0.30\cdot0.55\)\newline \(=0.165\) & \(0\) \\ \hline
		\(\mathrm{LL}\) & \(0.30\cdot0.45\)\newline \(=0.135\) & \(-6\) \\ \hline
	\end{tabularx}
\end{minipage}
\hfill
\begin{minipage}[t]{0.48\linewidth}
	\task{2}{Map the outcomes and build the PMF. Let \(Y\) be the payment adjustment in USD. The pairs \(\mathrm{TT}\), \(\mathrm{TL}\), \(\mathrm{LT}\), and \(\mathrm{LL}\) give \(8,0,0,-6\), respectively. Add the probabilities of the two pairs that give \(Y=0\).}
	
	\begin{tabularx}{\linewidth}{|c|Y|Y|}
		\hline
		\(y\) (USD) & Outcomes giving \(y\) & Calculation and mass \(p_Y(y)\) \\ \hline
		\(-6\) & \(\mathrm{LL}\) & \(0.135\) \\ \hline
		\(0\) & \(\mathrm{TL}\) and \(\mathrm{LT}\) & \(0.315+0.165\)\newline \(=0.480\) \\ \hline
		\(8\) & \(\mathrm{TT}\) & \(0.385\) \\ \hline
	\end{tabularx}
	
	\medskip
	The support is \(\{-6,0,8\}\) in USD.
	All masses are nonnegative and
	\(0.135+0.480+0.385=1\).
\end{minipage}

\task{3}{Find the probability of no payment adjustment. The PMF gives \(P(Y=0)=p_Y(0)=0.480\). The business has no adjustment in \(48\%\) of these invoice pairs.}

\clearpage
\begingroup
\setlength{\parskip}{2pt}
\stageheading{Guided practice}

\problemtitle{3}{Courier pay adjustment}{A}
Each delivery falls into exactly one category below. An early or on-time delivery adds \(2\) USD to the courier's base pay, a slightly late delivery leaves it unchanged, and a very late delivery deducts \(3\) USD.

\noindent
\noindent
\begin{minipage}[t]{0.48\linewidth}
	\task{1}{Read the four delivery outcomes and their supplied probabilities in the first two rows.}
	
	\begin{tabularx}{\linewidth}{|l|Y|Y|}
		\hline
		Complete outcome & Probability & \(A\) (USD) \\ \hline
		Early & \(0.18\) & \(2\) \\ \hline
		On time & \(0.47\) & \\ \hline
		Slightly late & \(0.23\) & \\ \hline
		Very late & \(0.12\) & \rule{0pt}{1.7em} \\ \hline
	\end{tabularx}
\end{minipage}
\hfill
\begin{minipage}[t]{0.48\linewidth}
	\task{2}{Define the pay adjustment \(A\), with units. Complete its mapping column on the left. Group outcomes by value to build the PMF below. Give the support and verify nonnegative masses and a sum of \(1\).}
	
	\begin{tabularx}{\linewidth}{|c|Y|Y|}
		\hline
		\(a\) (USD) & Outcomes giving \(a\) & Calculation and mass \(p_A(a)\) \\ \hline
		\(-3\) & & \rule{0pt}{1.55em} \\ \hline
		\(0\) & & \rule{0pt}{1.55em} \\ \hline
		\(2\) & & \rule{0pt}{1.55em} \\ \hline
	\end{tabularx}
\end{minipage}

\vspace{3mm}
\task{3}{Use the PMF to find \(P(A=2)\), the probability that the courier earns exactly \(2\) USD above base pay. Interpret the result.}


\problemtitle{4}{Delivery fees from two orders}{B}
Two customers independently keep (\(\mathrm{K}\)) or cancel (\(\mathrm{C}\)) their orders. Their cancellation probabilities are \(0.20\) and \(0.35\), respectively. The shop earns \(3\) USD per kept order and no fee per cancellation.

\noindent
\begin{minipage}[t]{0.48\linewidth}
	\task{1}{Complete the first two columns. List every keep/cancel outcome in customer order and multiply within each pair to calculate its probability.}
	
	\begin{tabularx}{\linewidth}{|l|Y|Y|}
		\hline
		Ordered outcome & Probability & \(F\) (USD) \\ \hline
		\(\mathrm{KK}\) & \(0.80\cdot0.65=0.52\) & \\ \hline
		& & \rule{0pt}{1.7em} \\ \hline
		& & \rule{0pt}{1.7em} \\ \hline
		& & \rule{0pt}{1.7em} \\ \hline
	\end{tabularx}
\end{minipage}
\hfill
\begin{minipage}[t]{0.48\linewidth}
	\task{2}{Define the total fee revenue \(F\), with units. Complete its mapping column on the left. Group outcomes by value to build the PMF below. Give the support and verify nonnegative masses and a sum of \(1\).}
	
	\begin{tabularx}{\linewidth}{|c|Y|Y|}
		\hline
		\(f\) (USD) & Outcomes giving \(f\) & Calculation and mass \(p_F(f)\) \\ \hline
		\(0\) & & \rule{0pt}{1.55em} \\ \hline
		\(3\) & & \rule{0pt}{1.55em} \\ \hline
		\(6\) & & \rule{0pt}{1.55em} \\ \hline
	\end{tabularx}
\end{minipage}

\task{3}{Use the PMF to find \(P(F=3)\), the probability that the shop earns exactly \(3\) USD in delivery fees. Interpret the result.}


\stageheading{Independent practice}
\vspace{-7mm}
\problemtitle{5}{Buying a carton of fruit}{A}
A food stall buys one carton of fruit. Exactly one supplier fills the order. The probabilities for the orchard, cooperative, wholesaler, and grocery are \(0.28\), \(0.19\), \(0.31\), and \(0.22\), respectively. Their prices for the carton are \(14\), \(18\), \(14\), and \(22\) USD, in the same order.

\task{1}{Record every supplier outcome and its given probability.}
\task{2}{Define \(C\), the purchase cost, with units. Map every outcome to its value of \(C\). List the support. Group outcomes by value to build the PMF, show your calculations, and verify nonnegative masses and a sum of \(1\).}
\task{3}{Use the PMF to find \(P(C=14)\), the probability that the stall pays exactly \(14\) USD for the carton. Interpret the result.}

\vspace{-1mm}
\problemtitle{6}{Demand and a supplier surcharge}{B}
\vspace{-1mm}
On a market day, a shop has high demand (\(\mathrm{H}\)) with probability \(0.30\) and regular demand (\(\mathrm{R}\)) otherwise. Independently, its supplier charges a surcharge (\(\mathrm{S}\)) with probability \(0.25\) and charges no surcharge (\(\mathrm{N}\)) otherwise. High demand adds \(6\) USD to profit relative to the shop's baseline, and the surcharge subtracts \(6\) USD. Regular demand and no surcharge each contribute \(0\) USD.

\task{1}{List every complete outcome in demand--surcharge order and multiply within each pair to calculate its probability.}
\task{2}{Define \(M\), the net change in profit, with units. Map every outcome to its value of \(M\). List the support. Group outcomes by value to build the PMF, show your calculations, and verify nonnegative masses and a sum of \(1\).}
\task{3}{Use the PMF to find \(P(M=0)\), the probability that the day's profit equals the baseline. Interpret the result.}

\vspace{-1mm}
\problemtitle{7}{Market stall rental credit}{A}
\vspace{-1mm}

A vendor is assigned exactly one market location. The standard stall fee is \(40\) USD. The table gives the fee the vendor actually pays at each location. The rental credit is the difference between the standard fee and the fee paid.

\begin{tabularx}{\linewidth}{@{}lYYYYY@{}}
    \toprule
    Location & Front gate & Main aisle & Side aisle & Rear row & Outdoor row \\
    \midrule
    Assignment probability & \(0.12\) & \(0.23\) & \(0.18\) & \(0.27\) & \(0.20\) \\
    Fee paid (USD) & \(40\) & \(40\) & \(30\) & \(30\) & \(20\) \\
    \bottomrule
\end{tabularx}

\task{1}{Record every market location outcome and its assignment probability.}
\task{2}{Define \(R\), the rental credit, with units. Subtract the fee paid from \(40\) USD to map each location to its credit. List the support. Group outcomes by value to build the PMF, show your calculations, and verify nonnegative masses and a sum of \(1\).}
\task{3}{Use the PMF to find \(P(R=10)\), the probability that the vendor receives exactly \(10\) USD in rental credit. Interpret the result.}

\vspace{-1mm}
\problemtitle{8}{An emergency restock trip}{B}
\vspace{-1mm}
Two stores independently request (\(\mathrm{R}\)) or do not request (\(\mathrm{N}\)) an emergency restock. The probabilities of a request are \(0.15\) for Store 1 and \(0.40\) for Store 2. If at least one requests a restock, the distributor makes a single trip at a cost of \(9\) USD plus \(3\) USD for each requesting store. With no requests, the cost is \(0\) USD.

\task{1}{List every complete request/no-request outcome in store order and multiply within each pair to calculate its probability.}
\task{2}{Define \(D\), the distributor's restock cost, with units. Map every outcome to its value of \(D\). List the support. Group outcomes by value to build the PMF, show your calculations, and verify nonnegative masses and a sum of \(1\).}
\task{3}{Use the PMF to find \(P(D=12)\), the probability that the distributor spends exactly \(12\) USD on the emergency restock. Interpret the result.}

\vspace{1.6in}\vfill

\ifprintsolutions
\clearpage
\RestoreClassroomHighlights
\input{graphs/G10_T2_bar_graphs.tex}
% Solution rows are indivisible; every child uses its own local linewidth.
% Compact height labels in the dense eight-outcome probability panels only.
\ExplSyntaxOn
\newcommand{\GTenDenseProbabilityGraph}[1]{%
  \group_begin:
  \pgfplotsset{every~axis~post/.append~style={
    execute~at~begin~axis={\tikzset{gtenbar~y~marker/.append~style={rotate=-90}}},
    yticklabel~style={rotate=90}}}
  #1\group_end:}
\ExplSyntaxOff
\newcommand{\GTenSolutionOutcomes}[3]{%
  \par\medskip\noindent\begin{minipage}{\linewidth}
  \begin{minipage}[t]{.40\linewidth}\vspace{0pt}\small
    \setlength{\tabcolsep}{3pt}#1
  \end{minipage}\hfill%
  \begin{minipage}[t]{.28\linewidth}\vspace{0pt}\raggedright#2\end{minipage}\hfill%
  \begin{minipage}[t]{.28\linewidth}\vspace{0pt}\raggedright#3\end{minipage}
  \end{minipage}\par\medskip}
\newcommand{\GTenSolutionPMF}[2]{%
  \par\medskip\noindent\begin{minipage}{\linewidth}
  \begin{minipage}[t]{.68\linewidth}\vspace{0pt}\small
    \setlength{\tabcolsep}{4pt}#1
  \end{minipage}\hfill%
  \begin{minipage}[t]{.30\linewidth}\vspace{0pt}\raggedright#2\end{minipage}
  \end{minipage}\par\medskip}

\section*{Solutions}

\GTenHLOutcome{Complete outcomes}\quad
\GTenHLValue{Variable values and units}\quad
\GTenHLProbability{Probabilities}

\problemtitle{3}{Courier pay adjustment}{A}
\task{1}{\GTenHLOutcome{Early}, \GTenHLOutcome{on time}, \GTenHLOutcome{slightly late}, and \GTenHLOutcome{very late} have probabilities \(\GTenHLProbability{0.18}\), \(\GTenHLProbability{0.47}\), \(\GTenHLProbability{0.23}\), and \(\GTenHLProbability{0.12}\), respectively.}

\GTenSolutionOutcomes{\begin{tabularx}{\linewidth}{|>{\raggedright\arraybackslash}p{.32\linewidth}|>{\raggedright\arraybackslash}p{.40\linewidth}|>{\raggedright\arraybackslash}X|}
\hline
\GTenHLOutcome{Complete} \GTenHLOutcome{outcome} & \GTenHLProbability{Probability} & \(\GTenHLValue{A}\) (USD) \\ \hline
\GTenHLOutcome{Early} & \(\GTenHLProbability{0.18}\) & \(\GTenHLValue{2}\) \\ \hline
\GTenHLOutcome{On} \GTenHLOutcome{time} & \(\GTenHLProbability{0.47}\) & \(\GTenHLValue{2}\) \\ \hline
\GTenHLOutcome{Slightly} \GTenHLOutcome{late} & \(\GTenHLProbability{0.23}\) & \(\GTenHLValue{0}\) \\ \hline
\GTenHLOutcome{Very} \GTenHLOutcome{late} & \(\GTenHLProbability{0.12}\) & \(\GTenHLValue{-3}\) \\ \hline
\end{tabularx}}{
\GTenOutcomeProbabilityGraph[width=\linewidth,height=4.7cm,rotate=35,xlabel={\GTenHLOutcome{Outcome}},xhighlight=GTenHLOutcome,ylabel={\GTenHLProbability{Probability}},yhighlight=GTenHLProbability,color=GTenHLProbability!32]{Early,On time,Slightly late,Very late}{0.18,0.47,0.23,0.12}
}{
\GTenOutcomeValueGraph[width=\linewidth,height=4.7cm,rotate=35,xlabel={\GTenHLOutcome{Outcome}},xhighlight=GTenHLOutcome,ystep=1,ylabel={$\GTenHLValue{A}(\GTenHLOutcome{\omega})$ \GTenHLValue{(USD)}},yhighlight=GTenHLValue,color=GTenHLValue!32]{Early,On time,Slightly late,Very late}{2,2,0,-3}
}

\task{2}{\(A\) measures the change from the courier's base pay for one delivery, in USD.
    \GTenHLOutcome{Early} and \GTenHLOutcome{on time} each give \(\GTenHLValue{2}\).
    \GTenHLOutcome{Slightly late} gives \(\GTenHLValue{0}\) and
    \GTenHLOutcome{very late} gives \(\GTenHLValue{-3}\). Add the early and on-time probabilities because these mutually exclusive outcomes both give \(A=2\).
    The PMF is \(p_A(-3)=\GTenHLProbability{0.12}\), \(p_A(0)=\GTenHLProbability{0.23}\), and
    \(p_A(2)=0.18+0.47=\GTenHLProbability{0.65}\).
    The support is \(\GTenHLValue{\{-3,0,2\}}\) in USD. All masses are nonnegative and
    \(0.12+0.23+0.65=1\).}

\GTenSolutionPMF{\begin{tabularx}{\linewidth}{|>{\raggedright\arraybackslash}p{.18\linewidth}|>{\raggedright\arraybackslash}p{.32\linewidth}|>{\raggedright\arraybackslash}X|}
\hline
\(\GTenHLValue{a}\) (USD) & Outcomes giving that value & Calculation and \GTenHLProbability{\(p_A(a)\)} \\ \hline
\(\GTenHLValue{-3}\) & \GTenHLOutcome{Very} \GTenHLOutcome{late} & \(\GTenHLProbability{0.12}\) \\ \hline
\(\GTenHLValue{0}\) & \GTenHLOutcome{Slightly} \GTenHLOutcome{late} & \(\GTenHLProbability{0.23}\) \\ \hline
\(\GTenHLValue{2}\) & \GTenHLOutcome{Early} \GTenHLOutcome{and} \GTenHLOutcome{on} \GTenHLOutcome{time} & \(0.18+0.47\)\newline \(\GTenHLProbability{=0.65}\) \\ \hline
\end{tabularx}}{
\GTenPMFGraph[width=\linewidth,height=4.7cm,xstep=1,xlabel={\GTenHLValue{$a$ (USD)}},ylabel={\GTenHLProbability{Probability}},xhighlight=GTenHLValue,yhighlight=GTenHLProbability,color=GTenHLProbability!32]{-3,0,2}{0.12,0.23,0.65}
}

\task{3}{From the PMF, \(P(A=2)=p_A(2)=\GTenHLProbability{0.65}\). The courier receives \(2\) USD above base pay with probability \(0.65\).}

\clearpage
\section*{Solutions (continued)}
\problemtitle{4}{Delivery fees from two orders}{B}
\task{1}{The probabilities of keeping the orders are \(1-0.20=0.80\) and \(1-0.35=0.65\). Multiply within each ordered outcome by independence.}


\GTenSolutionOutcomes{\begin{tabularx}{\linewidth}{|>{\raggedright\arraybackslash}p{.32\linewidth}|>{\raggedright\arraybackslash}p{.40\linewidth}|>{\raggedright\arraybackslash}X|}
\hline
\GTenHLOutcome{Complete} \GTenHLOutcome{outcome} & \GTenHLProbability{Probability} & \(\GTenHLValue{F}\) (USD) \\ \hline
\GTenHLOutcome{\(\mathrm{KK}\)} & \(0.80\cdot0.65\)\newline \(\GTenHLProbability{=0.52}\) & \(\GTenHLValue{6}\) \\ \hline
\GTenHLOutcome{\(\mathrm{KC}\)} & \(0.80\cdot0.35\)\newline \(\GTenHLProbability{=0.28}\) & \(\GTenHLValue{3}\) \\ \hline
\GTenHLOutcome{\(\mathrm{CK}\)} & \(0.20\cdot0.65\)\newline \(\GTenHLProbability{=0.13}\) & \(\GTenHLValue{3}\) \\ \hline
\GTenHLOutcome{\(\mathrm{CC}\)} & \(0.20\cdot0.35\)\newline \(\GTenHLProbability{=0.07}\) & \(\GTenHLValue{0}\) \\ \hline
\end{tabularx}}{
\GTenOutcomeProbabilityGraph[width=\linewidth,height=4.7cm,rotate=35,xlabel={\GTenHLOutcome{Outcome}},xhighlight=GTenHLOutcome,ylabel={\GTenHLProbability{Probability}},yhighlight=GTenHLProbability,color=GTenHLProbability!32]{KK,KC,CK,CC}{0.52,0.28,0.13,0.07}
}{
\GTenOutcomeValueGraph[width=\linewidth,height=4.7cm,rotate=35,xlabel={\GTenHLOutcome{Outcome}},xhighlight=GTenHLOutcome,ystep=3,ylabel={$\GTenHLValue{F}(\GTenHLOutcome{\omega})$ \GTenHLValue{(USD)}},yhighlight=GTenHLValue,color=GTenHLValue!32]{KK,KC,CK,CC}{6,3,3,0}
}

\task{2}{\(F\) is the total fee revenue in USD. Each kept order contributes \(3\) USD. The outcomes \(\mathrm{KK},\mathrm{KC},\mathrm{CK},\mathrm{CC}\) give \(\GTenHLValue{6},\GTenHLValue{3},\GTenHLValue{3},\GTenHLValue{0}\), respectively. Add across the mutually exclusive outcomes \(\mathrm{KC}\) and \(\mathrm{CK}\), which both give \(F=3\).
    The PMF is \(p_F(0)=\GTenHLProbability{0.07}\), \(p_F(3)=0.28+0.13=\GTenHLProbability{0.41}\), and
    \(p_F(6)=\GTenHLProbability{0.52}\).
    The support is \(\GTenHLValue{\{0,3,6\}}\) in USD. All masses are nonnegative and
    \(0.07+0.41+0.52=1\).}

\GTenSolutionPMF{\begin{tabularx}{\linewidth}{|>{\raggedright\arraybackslash}p{.18\linewidth}|>{\raggedright\arraybackslash}p{.32\linewidth}|>{\raggedright\arraybackslash}X|}
\hline
\(\GTenHLValue{f}\) (USD) & Outcomes giving that value & Calculation and \GTenHLProbability{\(p_F(f)\)} \\ \hline
\(\GTenHLValue{0}\) & \GTenHLOutcome{\(\mathrm{CC}\)} & \(\GTenHLProbability{0.07}\) \\ \hline
\(\GTenHLValue{3}\) & \GTenHLOutcome{\(\mathrm{KC}\) and \(\mathrm{CK}\)} & \(0.28+0.13\)\newline \(\GTenHLProbability{=0.41}\) \\ \hline
\(\GTenHLValue{6}\) & \GTenHLOutcome{\(\mathrm{KK}\)} & \(\GTenHLProbability{0.52}\) \\ \hline
\end{tabularx}}{
\GTenPMFGraph[width=\linewidth,height=4.7cm,xstep=3,xlabel={\GTenHLValue{$f$ (USD)}},ylabel={\GTenHLProbability{Probability}},xhighlight=GTenHLValue,yhighlight=GTenHLProbability,color=GTenHLProbability!32]{0,3,6}{0.07,0.41,0.52}
}

\task{3}{From the PMF, \(P(F=3)=p_F(3)=\GTenHLProbability{0.41}\). The shop earns exactly \(3\) USD in delivery fees with probability \(0.41\).}

\clearpage
\section*{Solutions (continued)}
\problemtitle{5}{Buying a carton of fruit}{A}
\task{1}{The \GTenHLOutcome{orchard}, \GTenHLOutcome{cooperative}, \GTenHLOutcome{wholesaler}, and \GTenHLOutcome{grocery} have probabilities \(\GTenHLProbability{0.28}\), \(\GTenHLProbability{0.19}\), \(\GTenHLProbability{0.31}\), and \(\GTenHLProbability{0.22}\), respectively.}

\GTenSolutionOutcomes{\begin{tabularx}{\linewidth}{|>{\raggedright\arraybackslash}p{.32\linewidth}|>{\raggedright\arraybackslash}p{.40\linewidth}|>{\raggedright\arraybackslash}X|}
\hline
\GTenHLOutcome{Complete} \GTenHLOutcome{outcome} & \GTenHLProbability{Probability} & \(\GTenHLValue{C}\) (USD) \\ \hline
\GTenHLOutcome{Orchard} & \(\GTenHLProbability{0.28}\) & \(\GTenHLValue{14}\) \\ \hline
\GTenHLOutcome{Cooperative} & \(\GTenHLProbability{0.19}\) & \(\GTenHLValue{18}\) \\ \hline
\GTenHLOutcome{Wholesaler} & \(\GTenHLProbability{0.31}\) & \(\GTenHLValue{14}\) \\ \hline
\GTenHLOutcome{Grocery} & \(\GTenHLProbability{0.22}\) & \(\GTenHLValue{22}\) \\ \hline
\end{tabularx}}{
\GTenOutcomeProbabilityGraph[width=\linewidth,height=4.7cm,rotate=35,xlabel={\GTenHLOutcome{Outcome}},xhighlight=GTenHLOutcome,ylabel={\GTenHLProbability{Probability}},yhighlight=GTenHLProbability,color=GTenHLProbability!32]{Orchard,Cooperative,Wholesaler,Grocery}{0.28,0.19,0.31,0.22}
}{
\GTenOutcomeValueGraph[width=\linewidth,height=4.7cm,rotate=35,xlabel={\GTenHLOutcome{Outcome}},xhighlight=GTenHLOutcome,ystep=5,ylabel={$\GTenHLValue{C}(\GTenHLOutcome{\omega})$ \GTenHLValue{(USD)}},yhighlight=GTenHLValue,color=GTenHLValue!32]{Orchard,Cooperative,Wholesaler,Grocery}{14,18,14,22}
}

\task{2}{\(C\) measures what the stall pays for one carton, in USD.
    The \GTenHLOutcome{orchard} and \GTenHLOutcome{wholesaler} each give \(\GTenHLValue{14}\).
    The \GTenHLOutcome{cooperative} gives \(\GTenHLValue{18}\) and the \GTenHLOutcome{grocery} gives \(\GTenHLValue{22}\). Add the orchard and wholesaler probabilities because only one supplier fills the order and both charge \(14\) USD.
    The PMF is \(p_C(14)=0.28+0.31=\GTenHLProbability{0.59}\), \(p_C(18)=\GTenHLProbability{0.19}\), and
    \(p_C(22)=\GTenHLProbability{0.22}\).
    The support is \(\GTenHLValue{\{14,18,22\}}\) in USD. All masses are nonnegative and
    \(0.59+0.19+0.22=1\).}

\GTenSolutionPMF{\begin{tabularx}{\linewidth}{|>{\raggedright\arraybackslash}p{.18\linewidth}|>{\raggedright\arraybackslash}p{.32\linewidth}|>{\raggedright\arraybackslash}X|}
\hline
\(\GTenHLValue{c}\) (USD) & Outcomes giving that value & Calculation and \GTenHLProbability{\(p_C(c)\)} \\ \hline
\(\GTenHLValue{14}\) & \GTenHLOutcome{Orchard} \GTenHLOutcome{and} \GTenHLOutcome{wholesaler} & \(0.28+0.31\)\newline \(\GTenHLProbability{=0.59}\) \\ \hline
\(\GTenHLValue{18}\) & \GTenHLOutcome{Cooperative} & \(\GTenHLProbability{0.19}\) \\ \hline
\(\GTenHLValue{22}\) & \GTenHLOutcome{Grocery} & \(\GTenHLProbability{0.22}\) \\ \hline
\end{tabularx}}{
\GTenPMFGraph[width=\linewidth,height=4.7cm,xstep=2,xlabel={\GTenHLValue{$c$ (USD)}},ylabel={\GTenHLProbability{Probability}},xhighlight=GTenHLValue,yhighlight=GTenHLProbability,color=GTenHLProbability!32]{14,18,22}{0.59,0.19,0.22}
}

\task{3}{From the PMF, \(P(C=14)=p_C(14)=\GTenHLProbability{0.59}\). The stall pays \(14\) USD for the carton with probability \(0.59\).}


\clearpage
\section*{Solutions (continued)}
\problemtitle{6}{Demand and a supplier surcharge}{B}
\task{1}{The regular-demand and no-surcharge probabilities are \(0.70\) and \(0.75\). Multiply within each ordered outcome by independence.}


\GTenSolutionOutcomes{\begin{tabularx}{\linewidth}{|>{\raggedright\arraybackslash}p{.32\linewidth}|>{\raggedright\arraybackslash}p{.40\linewidth}|>{\raggedright\arraybackslash}X|}
\hline
\GTenHLOutcome{Complete} \GTenHLOutcome{outcome} & \GTenHLProbability{Probability} & \(\GTenHLValue{M}\) (USD) \\ \hline
\GTenHLOutcome{\(\mathrm{HS}\)} & \(0.30\cdot0.25\)\newline \(\GTenHLProbability{=0.075}\) & \(\GTenHLValue{0}\) \\ \hline
\GTenHLOutcome{\(\mathrm{HN}\)} & \(0.30\cdot0.75\)\newline \(\GTenHLProbability{=0.225}\) & \(\GTenHLValue{6}\) \\ \hline
\GTenHLOutcome{\(\mathrm{RS}\)} & \(0.70\cdot0.25\)\newline \(\GTenHLProbability{=0.175}\) & \(\GTenHLValue{-6}\) \\ \hline
\GTenHLOutcome{\(\mathrm{RN}\)} & \(0.70\cdot0.75\)\newline \(\GTenHLProbability{=0.525}\) & \(\GTenHLValue{0}\) \\ \hline
\end{tabularx}}{
\GTenOutcomeProbabilityGraph[width=\linewidth,height=4.7cm,rotate=35,xlabel={\GTenHLOutcome{Outcome}},xhighlight=GTenHLOutcome,ylabel={\GTenHLProbability{Probability}},yhighlight=GTenHLProbability,color=GTenHLProbability!32]{HS,HN,RS,RN}{0.075,0.225,0.175,0.525}
}{
\GTenOutcomeValueGraph[width=\linewidth,height=4.7cm,rotate=35,xlabel={\GTenHLOutcome{Outcome}},xhighlight=GTenHLOutcome,ystep=3,ylabel={$\GTenHLValue{M}(\GTenHLOutcome{\omega})$ \GTenHLValue{(USD)}},yhighlight=GTenHLValue,color=GTenHLValue!32]{HS,HN,RS,RN}{0,6,-6,0}
}

\task{2}{\(M\) is the day's net change in profit from baseline, in USD. Adding the signed contributions gives \(6-6=\GTenHLValue{0}\) for \(\mathrm{HS}\), \(\GTenHLValue{6}\) for \(\mathrm{HN}\), \(\GTenHLValue{-6}\) for \(\mathrm{RS}\), and \(\GTenHLValue{0}\) for \(\mathrm{RN}\). Add across the mutually exclusive outcomes \(\mathrm{HS}\) and \(\mathrm{RN}\), which both give \(M=0\).
    The PMF is \(p_M(-6)=\GTenHLProbability{0.175}\), \(p_M(0)=0.075+0.525=\GTenHLProbability{0.600}\), and
    \(p_M(6)=\GTenHLProbability{0.225}\).
    The support is \(\GTenHLValue{\{-6,0,6\}}\) in USD. All masses are nonnegative and
    \(0.175+0.600+0.225=1\).}

\GTenSolutionPMF{\begin{tabularx}{\linewidth}{|>{\raggedright\arraybackslash}p{.18\linewidth}|>{\raggedright\arraybackslash}p{.32\linewidth}|>{\raggedright\arraybackslash}X|}
\hline
\(\GTenHLValue{m}\) (USD) & Outcomes giving that value & Calculation and \GTenHLProbability{\(p_M(m)\)} \\ \hline
\(\GTenHLValue{-6}\) & \GTenHLOutcome{\(\mathrm{RS}\)} & \(\GTenHLProbability{0.175}\) \\ \hline
\(\GTenHLValue{0}\) & \GTenHLOutcome{\(\mathrm{HS}\) and \(\mathrm{RN}\)} & \(0.075+0.525\)\newline \(\GTenHLProbability{=0.600}\) \\ \hline
\(\GTenHLValue{6}\) & \GTenHLOutcome{\(\mathrm{HN}\)} & \(\GTenHLProbability{0.225}\) \\ \hline
\end{tabularx}}{
\GTenPMFGraph[width=\linewidth,height=4.7cm,xstep=6,xlabel={\GTenHLValue{$m$ (USD)}},ylabel={\GTenHLProbability{Probability}},xhighlight=GTenHLValue,yhighlight=GTenHLProbability,color=GTenHLProbability!32]{-6,0,6}{0.175,0.600,0.225}
}

\task{3}{From the PMF, \(P(M=0)=p_M(0)=\GTenHLProbability{0.600}\). The shop's profit equals its baseline profit with probability \(0.600\).}

\clearpage
\section*{Solutions (continued)}
\problemtitle{7}{Market stall rental credit}{A}
\task{1}{\GTenHLOutcome{Front gate}, \GTenHLOutcome{main aisle}, \GTenHLOutcome{side aisle}, \GTenHLOutcome{rear row}, and \GTenHLOutcome{outdoor row} have probabilities \(\GTenHLProbability{0.12}\), \(\GTenHLProbability{0.23}\), \(\GTenHLProbability{0.18}\), \(\GTenHLProbability{0.27}\), and \(\GTenHLProbability{0.20}\), respectively.}

\GTenSolutionOutcomes{\begin{tabularx}{\linewidth}{|>{\raggedright\arraybackslash}p{.32\linewidth}|>{\raggedright\arraybackslash}p{.40\linewidth}|>{\raggedright\arraybackslash}X|}
\hline
\GTenHLOutcome{Complete} \GTenHLOutcome{outcome} & \GTenHLProbability{Probability} & \(\GTenHLValue{R}\) (USD) \\ \hline
\GTenHLOutcome{Front} \GTenHLOutcome{gate} & \(\GTenHLProbability{0.12}\) & \(\GTenHLValue{0}\) \\ \hline
\GTenHLOutcome{Main} \GTenHLOutcome{aisle} & \(\GTenHLProbability{0.23}\) & \(\GTenHLValue{0}\) \\ \hline
\GTenHLOutcome{Side} \GTenHLOutcome{aisle} & \(\GTenHLProbability{0.18}\) & \(\GTenHLValue{10}\) \\ \hline
\GTenHLOutcome{Rear} \GTenHLOutcome{row} & \(\GTenHLProbability{0.27}\) & \(\GTenHLValue{10}\) \\ \hline
\GTenHLOutcome{Outdoor} \GTenHLOutcome{row} & \(\GTenHLProbability{0.20}\) & \(\GTenHLValue{20}\) \\ \hline
\end{tabularx}}{
\GTenOutcomeProbabilityGraph[width=\linewidth,height=4.7cm,rotate=35,xlabel={\GTenHLOutcome{Outcome}},xhighlight=GTenHLOutcome,ylabel={\GTenHLProbability{Probability}},yhighlight=GTenHLProbability,color=GTenHLProbability!32]{Front gate,Main aisle,Side aisle,Rear row,Outdoor row}{0.12,0.23,0.18,0.27,0.20}
}{
\GTenOutcomeValueGraph[width=\linewidth,height=4.7cm,rotate=35,xlabel={\GTenHLOutcome{Outcome}},xhighlight=GTenHLOutcome,ystep=5,ylabel={$\GTenHLValue{R}(\GTenHLOutcome{\omega})$ \GTenHLValue{(USD)}},yhighlight=GTenHLValue,color=GTenHLValue!32]{Front gate,Main aisle,Side aisle,Rear row,Outdoor row}{0,0,10,10,20}
}

\task{2}{\(R\) measures the discount from the standard stall fee, in USD. Subtract the fee paid from \(40\).
    The \GTenHLOutcome{front gate} and \GTenHLOutcome{main aisle} each give \(40-40=\GTenHLValue{0}\).
    The \GTenHLOutcome{side aisle} and \GTenHLOutcome{rear row} each give \(40-30=\GTenHLValue{10}\).
    The \GTenHLOutcome{outdoor row} gives \(40-20=\GTenHLValue{20}\). Add within each group because the vendor receives exactly one location.
    The PMF is \(p_R(0)=0.12+0.23=\GTenHLProbability{0.35}\),
    \(p_R(10)=0.18+0.27=\GTenHLProbability{0.45}\), and \(p_R(20)=\GTenHLProbability{0.20}\).
    The support is \(\GTenHLValue{\{0,10,20\}}\) in USD. All masses are nonnegative and
    \(0.35+0.45+0.20=1\).}

\GTenSolutionPMF{\begin{tabularx}{\linewidth}{|>{\raggedright\arraybackslash}p{.18\linewidth}|>{\raggedright\arraybackslash}p{.32\linewidth}|>{\raggedright\arraybackslash}X|}
\hline
\(\GTenHLValue{r}\) (USD) & Outcomes giving that value & Calculation and \GTenHLProbability{\(p_R(r)\)} \\ \hline
\(\GTenHLValue{0}\) & \GTenHLOutcome{Front} \GTenHLOutcome{gate} \GTenHLOutcome{and} \GTenHLOutcome{main} \GTenHLOutcome{aisle} & \(0.12+0.23\)\newline \(\GTenHLProbability{=0.35}\) \\ \hline
\(\GTenHLValue{10}\) & \GTenHLOutcome{Side} \GTenHLOutcome{aisle} \GTenHLOutcome{and} \GTenHLOutcome{rear} \GTenHLOutcome{row} & \(0.18+0.27\)\newline \(\GTenHLProbability{=0.45}\) \\ \hline
\(\GTenHLValue{20}\) & \GTenHLOutcome{Outdoor} \GTenHLOutcome{row} & \(\GTenHLProbability{0.20}\) \\ \hline
\end{tabularx}}{
\GTenPMFGraph[width=\linewidth,height=4.7cm,xstep=10,xlabel={\GTenHLValue{$r$ (USD)}},ylabel={\GTenHLProbability{Probability}},xhighlight=GTenHLValue,yhighlight=GTenHLProbability,color=GTenHLProbability!32]{0,10,20}{0.35,0.45,0.20}
}

\task{3}{From the PMF, \(P(R=10)=p_R(10)=\GTenHLProbability{0.45}\). The vendor receives a \(10\) USD rental credit with probability \(0.45\).}

\clearpage
\section*{Solutions (continued)}
\problemtitle{8}{An emergency restock trip}{B}
\task{1}{The no-request probabilities are \(0.85\) and \(0.60\). Multiply within each ordered outcome by independence.}


\GTenSolutionOutcomes{\begin{tabularx}{\linewidth}{|>{\raggedright\arraybackslash}p{.32\linewidth}|>{\raggedright\arraybackslash}p{.40\linewidth}|>{\raggedright\arraybackslash}X|}
\hline
\GTenHLOutcome{Complete} \GTenHLOutcome{outcome} & \GTenHLProbability{Probability} & \(\GTenHLValue{D}\) (USD) \\ \hline
\GTenHLOutcome{\(\mathrm{RR}\)} & \(0.15\cdot0.40\)\newline \(\GTenHLProbability{=0.06}\) & \(\GTenHLValue{15}\) \\ \hline
\GTenHLOutcome{\(\mathrm{RN}\)} & \(0.15\cdot0.60\)\newline \(\GTenHLProbability{=0.09}\) & \(\GTenHLValue{12}\) \\ \hline
\GTenHLOutcome{\(\mathrm{NR}\)} & \(0.85\cdot0.40\)\newline \(\GTenHLProbability{=0.34}\) & \(\GTenHLValue{12}\) \\ \hline
\GTenHLOutcome{\(\mathrm{NN}\)} & \(0.85\cdot0.60\)\newline \(\GTenHLProbability{=0.51}\) & \(\GTenHLValue{0}\) \\ \hline
\end{tabularx}}{
\GTenOutcomeProbabilityGraph[width=\linewidth,height=4.7cm,rotate=35,xlabel={\GTenHLOutcome{Outcome}},xhighlight=GTenHLOutcome,ylabel={\GTenHLProbability{Probability}},yhighlight=GTenHLProbability,color=GTenHLProbability!32]{RR,RN,NR,NN}{0.06,0.09,0.34,0.51}
}{
\GTenOutcomeValueGraph[width=\linewidth,height=4.7cm,rotate=35,xlabel={\GTenHLOutcome{Outcome}},xhighlight=GTenHLOutcome,ystep=3,ylabel={$\GTenHLValue{D}(\GTenHLOutcome{\omega})$ \GTenHLValue{(USD)}},yhighlight=GTenHLValue,color=GTenHLValue!32]{RR,RN,NR,NN}{15,12,12,0}
}

\task{2}{\(D\) is the distributor's emergency restock cost in USD. For \(\mathrm{RR}\), the cost is \(9+2\cdot3=\GTenHLValue{15}\). For \(\mathrm{RN}\) or \(\mathrm{NR}\), it is \(9+3=\GTenHLValue{12}\). For \(\mathrm{NN}\), it is \(\GTenHLValue{0}\). Add across the mutually exclusive outcomes \(\mathrm{RN}\) and \(\mathrm{NR}\), which both give \(D=12\).
    The PMF is \(p_D(0)=\GTenHLProbability{0.51}\), \(p_D(12)=0.09+0.34=\GTenHLProbability{0.43}\), and
    \(p_D(15)=\GTenHLProbability{0.06}\).
    The support is \(\GTenHLValue{\{0,12,15\}}\) in USD. All masses are nonnegative and
    \(0.51+0.43+0.06=1\).}

\GTenSolutionPMF{\begin{tabularx}{\linewidth}{|>{\raggedright\arraybackslash}p{.18\linewidth}|>{\raggedright\arraybackslash}p{.32\linewidth}|>{\raggedright\arraybackslash}X|}
\hline
\(\GTenHLValue{d}\) (USD) & Outcomes giving that value & Calculation and \GTenHLProbability{\(p_D(d)\)} \\ \hline
\(\GTenHLValue{0}\) & \GTenHLOutcome{\(\mathrm{NN}\)} & \(\GTenHLProbability{0.51}\) \\ \hline
\(\GTenHLValue{12}\) & \GTenHLOutcome{\(\mathrm{RN}\) and \(\mathrm{NR}\)} & \(0.09+0.34\)\newline \(\GTenHLProbability{=0.43}\) \\ \hline
\(\GTenHLValue{15}\) & \GTenHLOutcome{\(\mathrm{RR}\)} & \(\GTenHLProbability{0.06}\) \\ \hline
\end{tabularx}}{
\GTenPMFGraph[width=\linewidth,height=4.7cm,xstep=3,xlabel={\GTenHLValue{$d$ (USD)}},ylabel={\GTenHLProbability{Probability}},xhighlight=GTenHLValue,yhighlight=GTenHLProbability,color=GTenHLProbability!32]{0,12,15}{0.51,0.43,0.06}
}

\task{3}{From the PMF, \(P(D=12)=p_D(12)=\GTenHLProbability{0.43}\). The distributor spends exactly \(12\) USD on the restock trip with probability \(0.43\).}
\fi

\end{document}
